\documentclass[11pt]{article}

\usepackage{fullpage}
\usepackage{amsmath}
\usepackage{amssymb}
\usepackage{amsfonts}
\usepackage{amsthm}
\usepackage{mathtools}
\usepackage{enumitem}
\usepackage[colorlinks,linkcolor=magenta,citecolor=blue]{hyperref}

\newcommand{\Rbar}{\bar R}
\newcommand{\Om}{\Omega}
\newcommand{\E}{\mathbb E}
\newcommand{\Var}{\operatorname{Var}}
\newcommand{\tr}{\operatorname{tr}}
\newcommand{\diag}{\operatorname{diag}}
\newcommand{\Unif}{\operatorname{Unif}}
\newcommand{\IG}{\operatorname{IG}}
\newcommand{\Gam}{\operatorname{Gamma}}
\newcommand{\Normal}{\mathcal N}

\newtheorem{lemma}{Lemma}

\title{Proof of the ``irrelevance of the null SNPs'' of the SER-relax model}
\author{}
\date{}

\begin{document}
\maketitle

\section{Two helpful lemmas}

\begin{lemma}
Let $R$ be a $J \times J$ positive semi-definite matrix, and let $x \in \mathbb{R}^J$ be a vector partitioned as $x = \begin{pmatrix} x_1 \\ x_2 \end{pmatrix}$, where $x_1$ is a scalar. Let $S_1$ be the Schur complement of the top-left scalar element $R_{11}$ in $R$, and let $S$ be a $J \times J$ matrix formed by padding $S_1$ with zeros such that the first row and column are zero, and the bottom-right $(J-1) \times (J-1)$ block is $S_1$. For any positive scalar $\zeta$, the following relationship holds:
\[
x^T (R + \zeta S)^{-1} x = \frac{1}{1 + \zeta} \left( x^T R^{-1} x \right) + \frac{\zeta}{1 + \zeta} \left( x_1^2 R_{11}^{-1} \right)
\]
\end{lemma}

\begin{proof}
Partition $R$ into block form as $R = \begin{pmatrix} R_{11} & R_{12} \\ R_{21} & R_{22} \end{pmatrix}$. The Schur complement of $R_{11}$ is $S_1 = R_{22} - R_{21} R_{11}^{-1} R_{12}$.

Let $\tilde{R} = R + \zeta S$. Written in block form, this is:
\[
\tilde{R} = \begin{pmatrix} R_{11} & R_{12} \\ R_{21} & R_{22} + \zeta S_1 \end{pmatrix}
\]

The Schur complement of $R_{11}$ in $\tilde{R}$ is given by:
\[
\tilde{S}_1 = (R_{22} + \zeta S_1) - R_{21} R_{11}^{-1} R_{12} = S_1 + \zeta S_1 = (1 + \zeta) S_1
\]

Using the standard block matrix inversion formula, the inverse of $R$ can be written as:
\[
R^{-1} = \begin{pmatrix} R_{11}^{-1} & 0 \\ 0 & 0 \end{pmatrix} + K S_1^{-1} K^T
\]
where $K = \begin{pmatrix} -R_{11}^{-1} R_{12} \\ I \end{pmatrix}$.

Applying the same inversion formula to $\tilde{R}$, and noting that $K$ remains unchanged, we obtain:
\[
\tilde{R}^{-1} = \begin{pmatrix} R_{11}^{-1} & 0 \\ 0 & 0 \end{pmatrix} + K \tilde{S}_1^{-1} K^T = \begin{pmatrix} R_{11}^{-1} & 0 \\ 0 & 0 \end{pmatrix} + \frac{1}{1 + \zeta} K S_1^{-1} K^T
\]

From the decomposition of $R^{-1}$, we know that $K S_1^{-1} K^T = R^{-1} - \begin{pmatrix} R_{11}^{-1} & 0 \\ 0 & 0 \end{pmatrix}$. Substituting this into the equation for $\tilde{R}^{-1}$ yields:
\[
\tilde{R}^{-1} = \begin{pmatrix} R_{11}^{-1} & 0 \\ 0 & 0 \end{pmatrix} + \frac{1}{1 + \zeta} \left( R^{-1} - \begin{pmatrix} R_{11}^{-1} & 0 \\ 0 & 0 \end{pmatrix} \right)
\]
\[
\tilde{R}^{-1} = \frac{1}{1 + \zeta} R^{-1} + \frac{\zeta}{1 + \zeta} \begin{pmatrix} R_{11}^{-1} & 0 \\ 0 & 0 \end{pmatrix}
\]

Finally, multiplying by $x^T$ on the left and $x$ on the right isolates $x_1^2 R_{11}^{-1}$ in the second term, resulting in:
\[
x^T (R + \zeta S)^{-1} x = \frac{1}{1 + \zeta} \left( x^T R^{-1} x \right) + \frac{\zeta}{1 + \zeta} \left( x_1^2 R_{11}^{-1} \right)
\]
which completes the proof.
\end{proof}

\vspace{1em}

\begin{lemma}
Let $R$, $S$, and $\zeta$ be defined as in Lemma 1. The determinant of the perturbed matrix $R + \zeta S$ satisfies the following relationship:
\[
\det(R + \zeta S) = (1 + \zeta)^{J-1} \det(R)
\]
\end{lemma}

\begin{proof}
Using the determinant formula for block matrices, the determinant of $R$ can be expressed as the product of the determinant of its top-left block and the determinant of its corresponding Schur complement:
\[
\det(R) = \det(R_{11}) \det(S_1)
\]

Applying the same block determinant formula to the perturbed matrix $\tilde{R} = R + \zeta S$, and recalling from Lemma 1 that its Schur complement is $\tilde{S}_1 = (1 + \zeta) S_1$, we have:
\[
\det(R + \zeta S) = \det(R_{11}) \det((1 + \zeta) S_1)
\]

Because $S_1$ is a matrix of dimension $(J-1) \times (J-1)$, factoring out the scalar $(1 + \zeta)$ from the determinant yields:
\[
\det((1 + \zeta) S_1) = (1 + \zeta)^{J-1} \det(S_1)
\]

Substituting this scalar property back into the determinant equation gives:
\[
\det(R + \zeta S) = (1 + \zeta)^{J-1} \det(R_{11}) \det(S_1)
\]

Since $\det(R_{11}) \det(S_1) = \det(R)$, we obtain the final result:
\[
\det(R + \zeta S) = (1 + \zeta)^{J-1} \det(R)
\]
which completes the proof.
\end{proof}


\section{Proof of the property}
Here we will show that the PIP from SER-relax does not depend on the off-diagonal parts of $\bar{R}$. Recall the SER-relax model:
\[
\begin{cases}
x  \mid z_j, \beta, \gamma=j
  \sim \mathcal{N} \left(\beta u_j(z_j),\; \dfrac{\bar{R}}{N}\right) \\
z_j \sim \mathcal{N}(\mu_j,\;\tau^2 S_j)
\end{cases}
\]

\noindent where $\mu_j = \bar{R}_{-j,j}\in \mathbb{R}^{J-1}$, \quad $S_j = \bar{R}_{-j,-j} - \mu_j\mu_j^T$ is the Schur complement of $\bar{R}_{jj}$, and $u_j(z_j)$ is the vector with 1 on the $j$-element and the rest is $z_j$.

\paragraph{Goal.} We want to calculate $p(\gamma = j \mid x,\beta)$ and $p(\gamma = j \mid x)$, and show that it only depends on $\bar{R}$ through
\[
q = x^T\bar{R}^{-1}x.
\]



\paragraph{Calculation for $j = 1$.}
\[
\begin{cases}
x = \begin{pmatrix} x_1 \\ x_{-1} \end{pmatrix} \mid z_1, \beta, \gamma=1
  \sim \mathcal{N} \left(\beta\begin{pmatrix}1\\z_1\end{pmatrix},\;
  \dfrac{\bar{R}}{N}\right) \\
z_1 \sim \mathcal{N}(\mu_1,\;\tau^2 S_1)
\end{cases}
\]

\noindent where $\mu_1 = \bar{R}_{-1,1}$, \quad $S_1 = \bar{R}_{-1,-1} - \mu_1\mu_1^T$.

\medskip
\noindent \textbf{Marginal \& Conditional:} Marginalize out $z_1$:

\[
x \mid \beta,\,\gamma=1 \;\sim\;
\mathcal{N}\!\left(
  \beta\begin{pmatrix}1\\\mu_1\end{pmatrix},\;
  \frac{1}{N}\begin{pmatrix}1 & \mu_1^T \\ \mu_1 & \bar{R}_{-1,-1}+N\beta^2\tau^2 S_1\end{pmatrix}
\right)
\]

\noindent Let
\[
\tilde{R}
= \begin{pmatrix} 1 & \mu_1^T \\ \mu_1 & \bar{R}_{-1,-1}+N\beta^2\tau^2 S_1 \end{pmatrix}
= \bar{R} + \begin{pmatrix} 0 & 0 \\ 0 & N\beta^2\tau^2 S_1 \end{pmatrix}.
\]
We have
\[
p(x \mid \beta,\,\gamma=1)
= \frac{N^{J/2}}{(2\pi)^{J/2}\,|\tilde{R}|^{1/2}}
  \exp\!\left(
    -\frac{N}{2}
    \left(x-\beta\begin{pmatrix}1\\\mu_1\end{pmatrix}\right)^{\!T}
    \tilde{R}^{-1}
    \left(x-\beta\begin{pmatrix}1\\\mu_1\end{pmatrix}\right)
  \right)
\]

\begin{enumerate}
\item Because $\begin{pmatrix}1\\\mu_1\end{pmatrix}$ is the first column of $\tilde{R}$,
\[
\tilde{R}^{-1}\begin{pmatrix}1\\\mu_1\end{pmatrix} = e_1,
\qquad
\begin{pmatrix}1\\\mu_1\end{pmatrix}^{\!T}
\tilde{R}^{-1}
\begin{pmatrix}1\\\mu_1\end{pmatrix} = 1.
\]

\item \textbf{Using Lemma 1} with $\zeta = N\beta^2\tau^2$ we have
\[
x^T\tilde{R}^{-1}x
= \frac{1}{1+N\beta^2\tau^2}\,x^T\bar{R}^{-1}x
  + \frac{N\beta^2\tau^2}{1+N\beta^2\tau^2}\,x_1^2.
\]


\item \textbf{Using Lemma 2:}
\[
\det(\tilde{R}) = (1+N\beta^2\tau^2)^{J-1}\,\det(\bar{R}).
\]
\end{enumerate}

\medskip
\noindent Combine all of them and let
\[
q = x^T\bar{R}^{-1}x,
\]
we have
\begin{align*}
p(x \mid \beta,\,\gamma=1)
&= \frac{N^{J/2}}{(2\pi)^{J/2}\,|\bar{R}|^{1/2}\,(1+N\beta^2\tau^2)^{(J-1)/2}}
\exp\!\left(-\frac{N}{2(1+N\beta^2\tau^2)}q\right) \\
&\quad\times
\exp\!\left(-\frac{N^2\beta^2\tau^2}{2(1+N\beta^2\tau^2)}x_1^2\right)
\exp\!\left(N\beta x_1 - \frac{N}{2}\beta^2\right).
\end{align*}

\noindent Similarly,
\begin{align}\label{eq:conditional-prob-x-given-beta-gamma}
p(x \mid \beta,\,\gamma=j)
&= \frac{N^{J/2}}{(2\pi)^{J/2}\,|\bar{R}|^{1/2}\,(1+N\beta^2\tau^2)^{(J-1)/2}}
\exp\!\left(-\frac{N}{2(1+N\beta^2\tau^2)}q\right) \nonumber \\
&\quad\times
\exp\!\left(-\frac{N^2\beta^2\tau^2}{2(1+N\beta^2\tau^2)}x_j^2\right)
\exp\!\left(N\beta x_j - \frac{N}{2}\beta^2\right).
\end{align}

% \paragraph{Conclusion for $p(\gamma=j\mid x,\beta)$.}
% Assuming a uniform prior on $\gamma$, the terms involving $|\bar{R}|$ and $q$ do not depend on $j$ when $\beta$ is fixed, so they cancel in the normalization over $j$. Thus
% \begin{align*}
% p(\gamma=j \mid x,\beta)
% &=
% \frac{
% \exp\!\left(-\frac{N^2\beta^2\tau^2}{2(1+N\beta^2\tau^2)}x_j^2 + N\beta x_j\right)}
% {\sum_{k=1}^J
% \exp\!\left(-\frac{N^2\beta^2\tau^2}{2(1+N\beta^2\tau^2)}x_k^2 + N\beta x_k\right)}.
% \end{align*}
% Therefore $p(\gamma=j\mid x,\beta)$ does not depend on $q=x^T\bar{R}^{-1}x$.

\paragraph{Conclusion for $p(\gamma=j\mid x)$.}
After integrating over $\beta$, the $q$-dependent term remains inside the integral. Dropping constants that are common to all $j$,
\begin{align*}
p(\gamma=j \mid x)
&\propto
\int
(1+N\beta^2\tau^2)^{-(J-1)/2}
\exp\!\left(-\frac{N}{2(1+N\beta^2\tau^2)}q\right) \\
&\quad\times
\exp\!\left(-\frac{N^2\beta^2\tau^2}{2(1+N\beta^2\tau^2)} x_j^2 \right) \times
\exp\!\left(N\beta x_j - \frac{N}{2}\beta^2\right)
p(\beta \mid 0,\sigma_0^2)\,d\beta .
\end{align*}
Now combine the last exponential factor with the Gaussian prior on $\beta$.
Let
\[
s_N^2 = \frac{1}{N+\sigma_0^{-2}},
\qquad
m_j = \frac{N}{N+\sigma_0^{-2}}x_j,
\]
so that $m_j \approx x_j$ and $s_N^2$ is small when $N$ is large. Completing
the square gives
\[
\exp\!\left(N\beta x_j - \frac{N}{2}\beta^2\right)
p(\beta \mid 0,\sigma_0^2)
=
C_j\,p(\beta \mid m_j,s_N^2),
\]
where
\[
C_j =
\frac{1}{\sigma_0\sqrt{N+\sigma_0^{-2}}}
\exp\!\left(
  \frac{\sigma_0^{2} N}{\sigma_0^{2} N+1}\frac{Nx_j^2}{2}
\right).
\]
Therefore the same expression can be written as
\begin{align*}
p(\gamma=j \mid x)
&\propto
\E_{\beta \sim \Normal(m_j,s_N^2)}
\Bigg[
(1+N\beta^2\tau^2)^{-(J-1)/2}
\exp\!\left(-\frac{N}{2(1+N\beta^2\tau^2)}q\right) \\
&\hspace{8em}\times
\exp\!\left(\left( \frac{N\sigma_0^2}{1+N\sigma_0^2} - \frac{N\beta^2\tau^2}{1+N\beta^2\tau^2} \right) \dfrac{N x_j^2}{2}\right)
\Bigg].
\end{align*}

\paragraph{Implication for PIP diffuseness.}
This expectation form gives a useful way to compare SER-relax to
SER-relax-prime when $\nu_0$ is small. Define
\[
z_j = \sqrt{N}x_j,\qquad
\rho = \frac{N\sigma_0^2}{1+N\sigma_0^2},
\qquad
r(a) = \frac{a}{1+a}.
\]
It is useful to change variables to
\[
U=\sqrt{N}\beta.
\]
Since
\[
\beta \mid x_j \sim \Normal(m_j,s_N^2),
\qquad
m_j = \rho x_j,\qquad s_N^2=\frac{1}{N+\sigma_0^{-2}},
\]
we have the exact induced distribution
\[
U \mid x_j \sim \Normal(\rho z_j,\rho).
\]
When $N\sigma_0^2$ is large, $\rho\approx 1$, so this is approximately
\[
U \mid x_j \sim \Normal(z_j,1).
\]
Thus the original SER-relax expression can be written as
\[
p(\gamma=j \mid x)
\propto
\E_{U \sim \Normal(\rho z_j,\rho)}
\left[
G(U^2)
\exp\left\{
\frac{1}{2}\left(\rho-r(\tau^2U^2)\right)z_j^2
\right\}
\right],
\]
where
\[
G(U^2)
=
(1+\tau^2U^2)^{-(J-1)/2}
\exp\left\{
-\frac{Nq}{2(1+\tau^2U^2)}
\right\}.
\]
This is the main distinction from SER-relax-prime. In the original model, the
relaxation scale inside the integral is $\tau^2U^2$, and the distribution of
$U$ is centered at the SNP-specific value $\rho z_j$. Thus a SNP with a large
$z_j$ is allowed to put mass on large $U^2$, whereas a null SNP with
$z_j\approx 0$ has $U$ centered near zero. To reach the same large $U^2$ as
the top SNP, a null SNP must pay a Gaussian tail cost under
$\Normal(\rho z_j,\rho)$.

A simple plug-in approximation replaces $U$ by its mean $\rho z_j$:
\[
\log p(\gamma=j\mid x)
\approx
\frac{1}{2}\left\{\rho-r(\tau^2\rho^2 z_j^2)\right\}z_j^2
-\frac{J-1}{2}\log(1+\tau^2\rho^2 z_j^2)
-\frac{Nq}{2(1+\tau^2\rho^2 z_j^2)}
+ \text{constant}.
\]
When $\rho\approx 1$, this becomes the simpler approximation with
$\tau^2z_j^2$. This plug-in formula is only an approximation; the exact
statement is the expectation over $U \sim \Normal(\rho z_j,\rho)$ above.

In SER-relax-prime, $\beta^2$ in the covariance is replaced by the fixed prior
scale $\sigma_0^2$. Consequently the determinant and $q$ terms are common
across $j$ and cancel in the normalization over SNPs. The resulting single-SNP
weight is
\[
\tilde p(\gamma=j \mid x)
\propto
\exp\left\{
\frac{1}{2}
\left[
\rho-r(N\sigma_0^2\tau^2)
\right]z_j^2
\right\}.
\]
Thus SER-relax-prime applies the same attenuation to the $z_j^2$ evidence for
every SNP:
\[
\text{SER-relax-prime attenuation: } r(N\sigma_0^2\tau^2).
\]
The exact log-weight gap between the top SNP $j_\star$ and another SNP $k$ in
SER-relax-prime is
\[
\log \tilde p(\gamma=j_\star\mid x)
-
\log \tilde p(\gamma=k\mid x)
=
\frac{1}{2}
\left\{\rho-r(N\sigma_0^2\tau^2)\right\}
(z_{j_\star}^2-z_k^2).
\]
If $\sigma_0^2$ is learned and is of the order of the largest observed
$x_j^2$, then
\[
N\sigma_0^2 \approx z_{j_\star}^2.
\]
Therefore the prime attenuation can be similar to the original model's
plug-in attenuation at the top SNP. This observation is important: the
difference between the two models is not necessarily that the top SNP is less
attenuated in original SER-relax. The difference is that SER-relax-prime
applies the top-SNP-scale attenuation globally to every SNP. If
$N\sigma_0^2\tau^2$ is large, then
\[
r(N\sigma_0^2\tau^2)\approx \rho\approx 1,
\]
and the coefficient of every $z_j^2$ is close to zero. The log-weight gaps
collapse, and the PIPs become diffuse.

In original SER-relax, the analogous relaxation is not global. It is integrated
against
\[
U\mid x_j \sim \Normal(\rho z_j,\rho).
\]
For the top SNP, this distribution is centered near $z_{j_\star}$, so the
integral can use large $U^2$ values and obtain LD-column flexibility. For a
null SNP, the distribution is centered near zero, so the same large relaxation
requires a tail event. Consequently null SNPs do not receive the same effective
relaxation as the top SNP. This SNP-specific integration can preserve the
ranking and keep the highest PIP on the SNP with the largest $Z$-score, even
when the learned $\sigma_0^2$ is of the order of $\max_j x_j^2$ and
$\nu_0$ is small.

\paragraph{Taylor expansion around the prime scale.}
The preceding comparison can be made more explicit by expanding the original
SER-relax expectation around the fixed scale used by SER-relax-prime. Let
\[
a=N\sigma_0^2,\qquad b=1+\tau^2 a,\qquad S=U^2,\qquad \Delta=S-a.
\]
Write the log of the non-Gaussian part of the original expectation as
\[
H_j(s)
=
-\frac{J-1}{2}\log(1+\tau^2s)
-\frac{Nq}{2(1+\tau^2s)}
+\frac{1}{2}\left\{\rho-r(\tau^2s)\right\}z_j^2.
\]
Then
\[
p(\gamma=j\mid x)
\propto
\E_{U\sim \Normal(\rho z_j,\rho)}
\left[\exp\{H_j(U^2)\}\right].
\]
A Taylor expansion around $s=a$ gives
\[
H_j(a+\Delta)
=
H_j(a)+A_j\Delta+\frac{1}{2}B_j\Delta^2+O(\Delta^3),
\]
where
\[
H_j(a)
=
-\frac{J-1}{2}\log b
-\frac{Nq}{2b}
+\frac{1}{2}\left\{\rho-r(\tau^2a)\right\}z_j^2,
\]
\[
A_j
=
\frac{\tau^2}{2}
\left\{
\frac{Nq-z_j^2}{b^2}
-\frac{J-1}{b}
\right\},
\]
and
\[
B_j
=
\frac{\tau^4}{2}
\left\{
\frac{J-1}{b^2}
-\frac{2(Nq-z_j^2)}{b^3}
\right\}.
\]
The zeroth-order term is
\[
\exp\{H_j(a)\}
\propto
\exp\left\{
\frac{1}{2}\left[
\rho-r(N\sigma_0^2\tau^2)
\right]z_j^2
\right\},
\]
because the determinant and $q$ terms in $H_j(a)$ do not depend on $j$. This is
exactly the SER-relax-prime weight. Thus SER-relax-prime can be viewed as the
zeroth-order approximation to original SER-relax obtained by setting
\[
U^2=N\beta^2=N\sigma_0^2
\]
for every SNP.

The first correction is SNP-specific because the distribution of $\Delta=U^2-a$
depends on $z_j$. Under $U\sim \Normal(\rho z_j,\rho)$,
\[
\E_j[\Delta]
=
\rho+\rho^2z_j^2-a,
\]
and
\[
\Var_j(\Delta)
=
2\rho^2+4\rho^3z_j^2.
\]
Keeping terms through second order gives the approximation
\[
\log p(\gamma=j\mid x)
\approx
H_j(a)
+A_j\E_j[\Delta]
+\frac{1}{2}B_j\E_j[\Delta^2]
+\frac{1}{2}A_j^2\Var_j(\Delta)
+\text{constant}.
\]
This expression shows where SER-relax can depart from SER-relax-prime. If
$j_\star$ is the strongest SNP and the learned prior variance satisfies
$a=N\sigma_0^2\approx z_{j_\star}^2$, then
\[
\E_{j_\star}[U^2]\approx a
\]
up to an $O(1)$ correction when $a$ is large. Therefore the prime
zeroth-order approximation is locally reasonable for the top SNP. For a null
SNP with $z_j\approx 0$, however,
\[
\E_j[U^2]\approx \rho\approx 1,
\qquad
\E_j[\Delta]\approx 1-a,
\]
so $U^2$ is not close to $a$ at all. The Taylor expansion around $a$ is then
not a good local approximation for that SNP. This is precisely the difference:
SER-relax-prime evaluates all SNPs at the same top-SNP-scale value
$U^2=a$, whereas original SER-relax evaluates each SNP by integrating over a
distribution centered at its own $\rho z_j$.

This also clarifies the role of the $G(U^2)$ term. At zeroth order,
$G(a)$ is common across SNPs and cancels, which is what happens in
SER-relax-prime. In original SER-relax, $G(U^2)$ remains inside the
SNP-specific expectation. A strong SNP, whose $U$ distribution is centered far
from zero, can place mass on large $U^2$ values and access the corresponding
LD-column flexibility. A null SNP, whose $U$ distribution is centered near
zero, cannot access the same part of $G(U^2)$ without paying a Gaussian tail
cost. Therefore even if $\sigma_0^2$ is learned to be close to the largest
$x_j^2$, the original model does not give every SNP the same effective
relaxation scale.

\paragraph{Suggested diagnostics.}
For each SNP, define the normalized SER-relax integrand in the $U$ scale,
\[
\pi_j(dU)
\propto
\phi(U;\rho z_j,\rho)\,
G(U^2)
\exp\left\{
\frac{1}{2}\left(\rho-r(\tau^2U^2)\right)z_j^2
\right\}\,dU,
\]
where $\phi(\cdot;\mu,v)$ is the $\Normal(\mu,v)$ density. A direct diagnostic
is to compute
\[
\E_{\pi_j}[\tau^2U^2]
\]
for each SNP and compare it with the prime value
\[
N\sigma_0^2\tau^2.
\]
The expected pattern is that top SNPs have larger effective relaxation, null
SNPs have much smaller effective relaxation, and SER-relax-prime uses one
common relaxation scale for all SNPs. A second diagnostic is to compare the
actual SER-relax PIPs with the plug-in approximation obtained by replacing
$U$ by $\rho z_j$ in the expectation above.

Therefore $p(\gamma=j\mid x)$ depends on $\bar{R}$ through $q=x^T\bar{R}^{-1}x$ in general. It does not depend on the off-diagonal elements of $\bar{R}$ separately once $q$ is fixed.


\subsection{Approximate irrelevance of the null SNPs}

Another way to understand the behavior of SER-relax is to condition on the
marginal statistic $x_j$. Write
\[
x=\begin{pmatrix}x_j\\x_{-j}\end{pmatrix},
\qquad
\bar R
=
\begin{pmatrix}
1 & \mu_j^T\\
\mu_j & \bar R_{-j,-j}
\end{pmatrix},
\qquad
S_j=\bar R_{-j,-j}-\mu_j\mu_j^T.
\]
In both susie-rss and SER-relax, the marginal distribution of $x_j$ is
\[
x_j\mid \beta,\gamma=j \sim \Normal(\beta,1/N).
\]
Therefore the posterior based only on $x_j$ is
\[
\beta\mid x_j,\gamma=j
\sim
\Normal(m_j,s_N^2),
\qquad
m_j=\rho x_j,\qquad
s_N^2=\frac{\rho}{N},
\qquad
\rho=\frac{N\sigma_0^2}{1+N\sigma_0^2}.
\]
In susie-rss, the conditional distribution of the remaining coordinates is
independent of $\beta$:
\[
p(x_{-j}\mid x_j,\beta,\gamma=j)
=
\Normal(x_j\mu_j,S_j/N).
\]
This is the exact irrelevance property: once $x_j$ is fixed, $x_{-j}$ does not
change the relative evidence for SNP $j$ through $\beta$.

In SER-relax, the conditional mean is still independent of $\beta$, but the
conditional covariance is inflated by $\beta^2$:
\[
p(x_{-j}\mid x_j,\beta,\gamma=j)
=
\Normal\left(
x_j\mu_j,
\left(\frac{1}{N}+\beta^2\tau^2\right)S_j
\right).
\]
Integrating over $\beta\mid x_j$ gives
\[
p(x_{-j}\mid x_j,\gamma=j)
=
\int
\Normal\left(
x_{-j}\mid x_j\mu_j,\,
\left(\frac{1}{N}+\beta^2\tau^2\right)S_j
\right)
p(\beta\mid x_j,\gamma=j)\,d\beta
\]
\[
=
\E_{\beta\sim \Normal(m_j,s_N^2)}
\left[
\Normal\left(
x_{-j}\mid x_j\mu_j,\,
\left(\frac{1}{N}+\beta^2\tau^2\right)S_j
\right)
\right].
\]
This is a mixture of Gaussians with the same mean and the same shape matrix
$S_j$, but with SNP-specific scalar variance inflation. Since
$\beta\mid x_j$ is concentrated near $m_j=\rho x_j$, a plug-in approximation gives
\[
p(x_{-j}\mid x_j,\gamma=j)
\approx
\Normal\left(
x_j\mu_j,\,
\left(\frac{1}{N}+m_j^2\tau^2\right)S_j
\right).
\]
When $\rho\approx 1$, this becomes
\[
p(x_{-j}\mid x_j,\gamma=j)
\approx
\Normal\left(
x_j\mu_j,\,
\left(\frac{1}{N}+x_j^2\tau^2\right)S_j
\right).
\]

This gives an approximate null-SNP irrelevance statement. If SNP $j$ is null,
then $x_j=O(N^{-1/2})$, and under $\beta\mid x_j$,
\[
\beta^2 = O_p(N^{-1}).
\]
Therefore the scalar covariance inflation satisfies
\[
\frac{1}{N}+\beta^2\tau^2
=
\frac{1}{N}\{1+O_p(\tau^2)\}.
\]
For moderate $\tau^2$, the SER-relax conditional distribution is close to the
susie-rss conditional distribution:
\[
p(x_{-j}\mid x_j,\gamma=j)
\approx
\Normal(x_j\mu_j,S_j/N).
\]
In contrast, for a strong SNP with $z_j^2$ large, the inflation
$1+N\beta^2\tau^2$ can be much larger than 1 under
$\beta\mid x_j,\gamma=j$, allowing the active SNP's LD column to relax. This
is the approximate version of irrelevance: null SNPs see little variance
inflation after conditioning on their small $x_j$, while strong SNPs can
receive substantial LD-column flexibility.

This argument also shows why using $p(\beta\mid x,\gamma=j)$ directly is less
clean. The full posterior $p(\beta\mid x,\gamma=j)$ already contains
$x_{-j}$ through the Gaussian likelihood above, so it is circular for proving
irrelevance of $x_{-j}$. The sharper statement is obtained by first
conditioning on $x_j$ and then asking how much the remaining conditional
distribution changes as $\beta$ varies over
$\beta\mid x_j,\gamma=j$.

\section{TO-DO}
Understand the relationship between 
\begin{enumerate}
\item PIP only depends on $x$ and the diagonal of $R$.
\item $p(x_{-j} | x_j, \beta, \gamma = j)$ does not depend on $\beta$.
\end{enumerate}
Look Yuxin's thesis and read the SI of susie-rss paper more carefully.

Why does having $\beta^2$ in the variance make sense? Of course it follows from the IW or Wi model, but intuitively why does it make sense?

Is there another model that arises from the conditional distribution:
\begin{equation}\label{eq:conditional-dist-Matthew}
p(x_{-j}\mid x_j, \beta, \gamma=j) = 
\Normal\left(
x_j \bar R_{j, -j} ,\,
\left(\frac{1}{N}+ \dfrac{x_j^2}{\nu_0 + 3}\right)S_j
\right)?
\end{equation}

\section{Some new models with better approximation to susie-love}
Recall our original SER model from what is called `susie-love':
\begin{align}\label{eq:likelihood-x-R-mean-var} 
\begin{cases}
x | R, \beta, \gamma \sim \mathcal{N}(\beta R_{\gamma}, R / N)\\
R \sim \mathcal{IW}(\nu_0 \bar R, \nu_0 + J + 1).
\end{cases}
\end{align}
Note that the latent matrix $R$ appears in both mean and variance of the likelihood in~\eqref{eq:likelihood-x-R-mean-var}. This model is very natural for unknown in-sample LD $R$, but is hard to make inferences. Last time we approximated $R$ in the variance to be $\bar R$, which makes the inference much easier but might eventually lead to the loss of the ``irrelevance of the null SNPs'' property. In principle, it makes sense that when $L = 1$, PIP of SNP j calculation should only depend on $x_j$ but not on other $x_k$ and $\bar R_{jk}$ (``irrelevance of the null SNPs'' property for SER).

I think this original model is actually feasible for inference. Consider the following equivalent way to represent it.

\subsection{Model A (equivalent to the original SER-love model)} Let's write down $p(x | R, \beta, \gamma = 1)$, other cases of $\gamma$ are similar. Let 
$$d_1 = R_{11} \in \mathbb{R}_{+}, \quad z_1 = \dfrac{R_{-1, 1}}{d_1} \in \mathbb{R}^{J-1},$$
and 
$$S_1 = R_{-1, -1} - d_1 z_1 z_1^\top $$
be the Schur's complement of $R_{11}$ in $R$. Hence, the latent $R$ is represented as
\begin{equation}\label{eq:R-block}
R = \begin{pmatrix}
d_1 & d_1 z_1^\top \\
d_1 z_1 & S_1 + d_1 z_1 z_1^\top
\end{pmatrix}.
\end{equation}
Similarly, denote
$$\bar S_1 = \bar R_{-1, -1} - \bar R_{-1, 1} \bar R_{-1, 1}^\top $$
be the Schur's complement of $\bar R_{11}$ in $\bar R$, and $\mu_1 =  \bar R_{-1, 1}$, we can write
\begin{equation}\label{eq:R-bar-block}
\bar R = \begin{pmatrix}
1 &\mu_1^\top \\
\mu_1 & \bar S_1 + \mu_1 \mu_1^\top
\end{pmatrix}.
\end{equation}

Properties of IW distribution (together with the fact that $\bar R_{11} = 1$)
imply, using the shape--scale parameterization for the inverse-gamma
distribution,
\begin{equation}
\begin{cases}
d_1 \sim \IG\left(\dfrac{\nu_0 + 2}{2}, \dfrac{\nu_0}{2}\right) \\
z_1 | S_1 \sim \mathcal{N}\left(\bar R_{-1, 1}, \dfrac{S_1}{\nu_0}\right)\\
S_1 \sim \mathcal{IW}(\nu_0 \bar S_1, \nu_0 + J + 1)
\end{cases}
\end{equation}
The original SER version of Susie-love (with $\gamma = 1$; similar for other $\gamma$) becomes:
\begin{align}\label{eq:SER-model-A} 
\begin{cases}
x | \beta, \gamma = 1, z_1, S_1 \sim \mathcal{N}\left(\beta 
\begin{pmatrix}
d_1  \\
d_1 z_1
\end{pmatrix}, 
\dfrac{1}{N} \begin{pmatrix}
d_1 & d_1 z_1^\top \\
d_1 z_1 & S_1 + d_1 z_1 z_1^\top
\end{pmatrix}\right)\\
d_1 \sim \IG\left(\dfrac{\nu_0 + 2}{2}, \dfrac{\nu_0}{2}\right) \\
z_1 | S_1 \sim \mathcal{N}\left(\mu_1, \dfrac{S_1}{\nu_0}\right)\\
S_1 \sim \mathcal{IW}(\nu_0 \bar S_1, \nu_0 + J + 1)\\
\beta \sim \mathcal{N}(0, \sigma_0^2)
\end{cases}
\end{align}

\paragraph{Deriving $p(x\mid \gamma=1)$.}
Write
\[
x=\begin{pmatrix}x_1\\x_{-1}\end{pmatrix}.
\]
\paragraph{Step 1: likelihood conditional on all latent variables.}
From the block normal distribution in~\eqref{eq:SER-model-A},
\[
p(x\mid \gamma=1,\beta,d_1,z_1,S_1)
=
\Normal_J\left(
x\mid
\beta
\begin{pmatrix}
d_1\\ d_1z_1
\end{pmatrix},
\frac{1}{N}
\begin{pmatrix}
d_1 & d_1z_1^\top\\
d_1z_1 & S_1+d_1z_1z_1^\top
\end{pmatrix}
\right).
\]
Equivalently, using the conditional distribution of a block normal,
\[
p(x\mid \gamma=1,\beta,d_1,z_1,S_1)
=
\Normal\left(x_1\mid \beta d_1,\frac{d_1}{N}\right)
\Normal_{J-1}\left(x_{-1}\mid x_1z_1,\frac{S_1}{N}\right).
\]
The key simplification is that, \textbf{after conditioning on $x_1$, the conditional
mean of $x_{-1}$ is $x_1z_1$ and no longer contains $\beta$ or $d_1$} (relevant to our earlier discussion).

\paragraph{Step 2: marginalize over $z_1\mid S_1$.}
Since
\[
z_1\mid S_1\sim \Normal_{J-1}\left(\mu_1,\frac{S_1}{\nu_0}\right)
\]
and
\[
x_{-1}\mid x_1,z_1,S_1
\sim
\Normal_{J-1}\left(x_1z_1,\frac{S_1}{N}\right),
\]
we get
\begin{equation}\label{eq:conditional-dist-susie-love}
x_{-1}\mid x_1,S_1,\gamma=1
\sim
\Normal_{J-1}\left(
x_1\mu_1,
\left(\frac{1}{N}+\frac{x_1^2}{\nu_0}\right)S_1
\right).
\end{equation}
Therefore
\[
p(x\mid \gamma=1,\beta,d_1,S_1)
=
\Normal\left(x_1\mid \beta d_1,\frac{d_1}{N}\right)
\Normal_{J-1}\left(
x_{-1}\mid x_1\mu_1,
\left(\frac{1}{N}+\frac{x_1^2}{\nu_0}\right)S_1
\right).
\]
What I found so interesting is that the conditional distribution $p(x_{-1}\mid x_1,S_1,\gamma=1)$ resembles what was suggested by Matthew last time (cf.~\eqref{eq:conditional-dist-Matthew}).

\paragraph{Step 3: marginalize over $S_1$.}
Because $S_1\sim \mathcal{IW}_{J-1}(\nu_0\bar S_1,\nu_0+J+1)$, integrating
$S_1$ gives a multivariate $t$ distribution with $\nu_0+3$ degrees of freedom:
\[
x_{-1}\mid x_1,\gamma=1
\sim
t_{J-1,\nu_0+3}\left(
x_1\mu_1,\,
\frac{\nu_0}{\nu_0+3}\left(\frac{1}{N}+\frac{x_1^2}{\nu_0}\right)\bar S_1
\right),
\]
where the second argument is the multivariate $t$ scale matrix. Thus
\[
p(x\mid \gamma=1,\beta,d_1)
=
\Normal\left(x_1\mid \beta d_1,\frac{d_1}{N}\right)
t_{J-1,\nu_0+3}\left(
x_{-1}\mid x_1\mu_1,\,
\frac{\nu_0}{\nu_0+3}\left(\frac{1}{N}+\frac{x_1^2}{\nu_0}\right)\bar S_1
\right).
\]
Equivalently, the $t$ factor is proportional to
\[
|\bar S_1|^{-1/2}
\left(\frac{1}{N}+\frac{x_1^2}{\nu_0}\right)^{-(J-1)/2}
\left[
1+
\frac{
(x_{-1}-x_1\mu_1)^\top(\nu_0\bar S_1)^{-1}
(x_{-1}-x_1\mu_1)
}{\frac{1}{N}+\frac{x_1^2}{\nu_0}}
\right]^{-(\nu_0+J+2)/2}.
\]
Using the block inverse formula (notice the block form of $\bar R$~\eqref{eq:R-bar-block}),
\[
x^\top\bar R^{-1}x
=
x_1^2+
(x_{-1}-x_1\mu_1)^\top\bar S_1^{-1}
(x_{-1}-x_1\mu_1),
\]
so
\[
1+
\frac{
(x_{-1}-x_1\mu_1)^\top(\nu_0\bar S_1)^{-1}
(x_{-1}-x_1\mu_1)
}{\frac{1}{N}+\frac{x_1^2}{\nu_0}}
=
\frac{x^\top\bar R^{-1}x+\nu_0/N}{x_1^2+\nu_0/N}.
\]
Therefore
\begin{equation}\label{eq:T_1}
T_1(x)
\propto
|\bar S_1|^{-1/2}
\left(x_1^2+\frac{\nu_0}{N}\right)^{(\nu_0+3)/2}
\left(x^\top\bar R^{-1}x+\frac{\nu_0}{N}\right)^{-(\nu_0+J+2)/2}.
\end{equation}
When comparing different values of $\gamma$, the final factor is common across
SNPs. Also, since $\bar R_{11}=1$, $|\bar R|=|\bar S_1|$; analogously,
$|\bar S_j|=|\bar R|$ for any SNP $j$ with $\bar R_{jj}=1$.



\paragraph{Step 4: marginalize over $\beta$ and $d_1$.}
The preceding display separates the problem into an $x_{-1}\mid x_1$ factor
and an $x_1$ factor. Define
\[
T_1(x)=
t_{J-1,\nu_0+3}\left(
x_{-1}\mid x_1\mu_1,\,
\frac{\nu_0}{\nu_0+3}\left(\frac{1}{N}+\frac{x_1^2}{\nu_0}\right)\bar S_1
\right).
\]
If we integrate over $\beta$, then
\[
x_1\mid d_1,\gamma=1
\sim
\Normal\left(0,\frac{d_1}{N}+d_1^2\sigma_0^2\right),
\]
and hence (finally integrating over $d_1$):

\[
\boxed{
p(x\mid \gamma=1)
=
T_1(x)
\int_0^\infty
\Normal\left(x_1\mid 0,\frac{d_1}{N}+d_1^2\sigma_0^2\right)
\IG\left(d_1;\frac{\nu_0+2}{2},\frac{\nu_0}{2}\right)\,d d_1
}.
\]
This is a one-dimensional integral over $d_1>0$. I do not see a standard
closed form for it because the variance is $d_1/N+d_1^2\sigma_0^2$. Replacing $T_1$ with its expression in~\eqref{eq:T_1}, 

\[
\begin{aligned}
p(x\mid \gamma=1)
\propto\;& |\bar R|^{-1/2}
\left(x_1^2+\frac{\nu_0}{N}\right)^{(\nu_0+3)/2}
\left(x^\top\bar R^{-1}x+\frac{\nu_0}{N}\right)^{-(\nu_0+J+2)/2} \\
&\quad\times
\int_0^\infty
\Normal\left(x_1\mid 0,\frac{d_1}{N}+d_1^2\sigma_0^2\right)\,d \IG(d_1)
\end{aligned}
\]
Since $|\bar R|$ and $x^\top\bar R^{-1}x$ do not depend on which SNP is chosen,
both factors cancel in the PIP normalization. Also,
\[
\left(x_j^2+\frac{\nu_0}{N}\right)^{(\nu_0+3)/2}
=
\left(\frac{\nu_0}{N}\right)^{(\nu_0+3)/2}
\left(1+\frac{Nx_j^2}{\nu_0}\right)^{(\nu_0+3)/2},
\]
and the first factor is common across SNPs. Hence,
\[
\boxed{
p(\gamma=j\mid x)
\propto
\left(1+\frac{Nx_j^2}{\nu_0}\right)^{(\nu_0+3)/2}
\int_0^\infty
\Normal\left(x_j\mid 0,\frac{d_j}{N}+d_j^2\sigma_0^2\right)\,d \IG(d_j)
}.
\]
\textbf{Conclusion.} In the original SER-love model, the marginal likelihood
depends on $x^\top\bar R^{-1}x$, but this dependence is common across SNPs and
cancels when computing the PIP. The resulting PIP weight depends on SNP $j$
only through $x_j$.

\subsection{Model B} 
Simplify $d_1 \equiv 1$ to avoid intractable integral (and we know the diagonal of the in-sample LD matrix is 1 anyway), we have the following model (model B):
\begin{align}\label{eq:SER-model-B} 
\begin{cases}
x | \beta, \gamma = 1, z_1, S_1 \sim \mathcal{N}\left(\beta 
\begin{pmatrix}
1  \\
z_1
\end{pmatrix}, 
\dfrac{1}{N} \begin{pmatrix}
1 & z_1^\top \\
z_1 & S_1 + z_1 z_1^\top
\end{pmatrix}\right)\\
z_1 | S_1 \sim \mathcal{N}\left(\mu_1, \dfrac{S_1}{\nu_0}\right)\\
S_1 \sim \mathcal{IW}(\nu_0 \bar S_1, \nu_0 + J + 1)\\
\beta \sim \mathcal{N}(0, \sigma_0^2)
\end{cases}
\end{align}

This model leads to
\[
\boxed{
p(x\mid\gamma=1)
\propto
|\bar R|^{-1/2}
\left(1+\frac{Nx_j^2}{\nu_0}\right)^{(\nu_0+3)/2}
\left(x^\top\bar R^{-1}x+\frac{\nu_0}{N}\right)^{-(\nu_0+J+2)/2}
\Normal\left(x_1\mid 0,\frac{1}{N} + \sigma_0^2\right)
}.
\]
Since $|\bar R|$ and $x^\top\bar R^{-1}x$ do not depend on which SNP is chosen,
the PIP weights simplify further. Dropping the common factors gives
\begin{equation}\label{eq:PIP-model-B}
\boxed{
p(\gamma=j \mid x)
\propto
\left(1+\frac{Nx_j^2}{\nu_0}\right)^{(\nu_0+3)/2}
\Normal\left(x_j \mid 0,\frac{1}{N} + \sigma_0^2\right)
}.
\end{equation}

\paragraph{Compare to SER-RSS's PIP.} Recall SER-RSS's PIP:
\begin{equation}\label{eq:PIP-susie-RSS}
\tilde{p}(\gamma=j\mid x)
\propto \exp\left(\dfrac{N\sigma_0^{2}}{N\sigma_0^{2} + 1} \dfrac{N x_j^2}{2} \right)
\end{equation}
We now show that~\eqref{eq:PIP-model-B} tends to~\eqref{eq:PIP-susie-RSS}
as $\nu_0\to\infty$. Since
\[
\left(1+\frac{Nx_j^2}{\nu_0}\right)^{(\nu_0+3)/2} = \left(1+ \dfrac{2}{\nu_0 + 3} \times \dfrac{\nu_0 + 3}{\nu_0} \frac{Nx_j^2}{2}\right)^{(\nu_0+3)/2} 
\asymp \exp\left\{\frac{\nu_0 + 3}{\nu_0} \dfrac{Nx_j^2}{2}\right\}
\]
Also,
\[
\Normal\left(x_j\mid 0,\frac{1}{N}+\sigma_0^2\right)
\propto
\exp\left\{
-\frac{x_j^2}{2(1/N+\sigma_0^2)}
\right\}
=
\exp\left\{
-\frac{1}{1+N\sigma_0^2} \frac{N x_j^2}{2}
\right\}.
\]
Therefore the limiting weight in~\eqref{eq:PIP-model-B} is asymptotically (as $\nu_0$ gets large)
\[
\exp\left\{\left(\frac{\nu_0 + 3}{\nu_0} - \dfrac{1}{N\sigma_0^2 + 1}\right) \dfrac{Nx_j^2}{2}\right\} = \exp\left\{\left(\dfrac{N\sigma_0^2}{N\sigma_0^2 + 1} - \dfrac{1}{\nu_0 + 3}\right) \dfrac{Nx_j^2}{2}\right\} \approx \exp\left\{\dfrac{N\sigma_0^2}{N\sigma_0^2 + 1} \dfrac{Nx_j^2}{2}\right\}
\]
which is exactly the SER-RSS PIP weight in~\eqref{eq:PIP-susie-RSS}. Hence, for finite $\nu_0$, PIP under SER-love model (B) is slightly more diffused than susie-rss. 

\paragraph{Inference under model B~\eqref{eq:SER-model-B}.}
For a general SNP $j$, write
\[
\mu_j=\bar R_{-j,j},
\qquad
\bar S_j=\bar R_{-j,-j}-\mu_j\mu_j^\top,
\qquad
\kappa_j=\nu_0+Nx_j^2.
\]
Under model B, conditional on $\gamma=j$, the likelihood factorizes as
\[
p(x\mid \beta,z_j,S_j,\gamma=j)
=
\Normal\left(x_j\mid \beta,\frac{1}{N}\right)
\Normal_{J-1}\left(x_{-j}\mid x_jz_j,\frac{S_j}{N}\right).
\]
Thus $\beta$ is conditionally independent of $(z_j,S_j)$ given $x$ and
$\gamma=j$, and
\[
p(\beta,z_j,S_j\mid x,\gamma=j)
=
p(\beta\mid x_j,\gamma=j)\,
p(z_j,S_j\mid x,\gamma=j).
\]

$\bullet$ The posterior for $\beta$ is the usual normal-normal update:
\[
\boxed{
\beta\mid x,\gamma=j
\sim
\Normal\left(
\frac{N}{N+\sigma_0^{-2}}x_j,\,
\frac{1}{N+\sigma_0^{-2}}
\right).
}
\]

$\bullet$ For the LD-column variables, first condition on $S_j$. Combining
\[
z_j\mid S_j\sim \Normal_{J-1}\left(\mu_j,\frac{S_j}{\nu_0}\right),
\qquad
x_{-j}\mid x_j,z_j,S_j,\gamma=j
\sim
\Normal_{J-1}\left(x_jz_j,\frac{S_j}{N}\right)
\]
gives
\[
\boxed{
z_j\mid S_j,x,\gamma=j
\sim
\Normal_{J-1}\left(
\frac{\nu_0\mu_j+Nx_jx_{-j}}{\nu_0+Nx_j^2},\,
\frac{S_j}{\nu_0+Nx_j^2}
\right).
}
\]

$\bullet$ Integrating out $z_j$ gives
\[
x_{-j}\mid x_j,S_j,\gamma=j
\sim
\Normal_{J-1}\left(
x_j\mu_j,\,
\left(\frac{1}{N}+\frac{x_j^2}{\nu_0}\right)S_j
\right).
\]
Therefore the posterior for $S_j$ is inverse-Wishart:
\[
\boxed{
S_j\mid x,\gamma=j
\sim
\mathcal{IW}\left(
\nu_0\bar S_j+
\frac{N\nu_0}{\nu_0+Nx_j^2}
(x_{-j}-x_j\mu_j)(x_{-j}-x_j\mu_j)^\top,\,
\nu_0+J+2
\right).
}
\]
Equivalently, the joint posterior of $(z_j,S_j)$ can be represented by first
drawing $S_j$ from this inverse-Wishart distribution and then drawing $z_j$
from the conditional normal distribution above. Finally, the full posterior
over $\gamma$ uses the PIP weights in~\eqref{eq:PIP-model-B}, normalized over
$j=1,\ldots,J$.

\paragraph{Closed-form ELBO.}
Because the conditional posteriors above are exact, the optimized ELBO for
model B is the exact marginal log likelihood. Let
\[
V_\beta=\frac{1}{N}+\sigma_0^2,
\qquad
q=x^\top\bar R^{-1}x.
\]
For a uniform prior on $\gamma$, the optimized ELBO is
\[
\mathcal L(\nu_0,\sigma_0^2)
=
\log\left\{
\frac{1}{J}\sum_{j=1}^J \exp(\ell_j)
\right\},
\]
where the component log marginal likelihood can be written as
\begin{align*}
\ell_j
=\;&
-\frac{1}{2}\log(2\pi V_\beta)
-\frac{x_j^2}{2V_\beta}
+\log\Gamma\left(\frac{\nu_0+J+2}{2}\right)
-\log\Gamma\left(\frac{\nu_0+3}{2}\right) \\
&\quad
-\frac{J-1}{2}\log\pi
-\frac{1}{2}\log|\bar S_j|
+\frac{\nu_0+3}{2}\log\left(x_j^2+\frac{\nu_0}{N}\right)
-\frac{\nu_0+J+2}{2}\log\left(q+\frac{\nu_0}{N}\right).
\end{align*}
This is equivalent to the marginal likelihood derived above, but unlike the
PIP formula it keeps the terms that are common across SNPs. Those common terms
do not affect $\pi_j=p(\gamma=j\mid x)$ for fixed hyperparameters, but they do
matter when learning $\nu_0$.

\paragraph{Empirical Bayes learning of $\nu_0$ and $\sigma_0^2$.}
The empirical Bayes estimates can be obtained by maximizing the closed-form
objective
\[
(\hat\nu_0,\hat\sigma_0^2)
=
\arg\max_{\nu_0>0,\;\sigma_0^2\ge 0}
\mathcal L(\nu_0,\sigma_0^2).
\]
After choosing $(\hat\nu_0,\hat\sigma_0^2)$, the posterior inclusion
probabilities are
\[
\pi_j
=
\frac{\exp(\ell_j)}
{\sum_{k=1}^J \exp(\ell_k)}
\bigg|_{\nu_0=\hat\nu_0,\;\sigma_0^2=\hat\sigma_0^2}.
\]
Equivalently, for computing PIPs at fixed hyperparameters one may use the
simplified weights in~\eqref{eq:PIP-model-B}; for empirical Bayes optimization
one should use $\mathcal L(\nu_0,\sigma_0^2)$ above, because it retains the
$\nu_0$-dependent normalizing constants and the common
$\log(q+\nu_0/N)$ term.

\paragraph{EM view of empirical Bayes.}
Another way to use the fact that the posterior is exact is to write the usual
ELBO identity. Let
\[
\theta=(\gamma,\beta,z_\gamma,S_\gamma),
\qquad
q(\theta)=p(\theta\mid x;\nu_0^{\mathrm{old}},\sigma_{0,\mathrm{old}}^2).
\]
Strictly speaking,
\[
\log p(x;\nu_0,\sigma_0^2)
\ge
\E_q\{\log p(x\mid\theta)\}
-\operatorname{KL}\{q(\theta)\,\|\,p(\theta;\nu_0,\sigma_0^2)\},
\]
with equality when $q$ is the posterior under the same hyperparameters. In an
EM algorithm, the E-step fixes $q$ at the current hyperparameters and the
M-step maximizes
\[
Q(\nu_0,\sigma_0^2)
=
\E_q\{\log p(\theta;\nu_0,\sigma_0^2)\},
\]
because the likelihood term does not involve $\nu_0$ or $\sigma_0^2$ in model B.

This gives a closed-form update for $\sigma_0^2$. If
\[
\pi_j=q(\gamma=j),
\qquad
m_{\beta j}=\frac{N}{N+\sigma_{0,\mathrm{old}}^{-2}}x_j,
\qquad
v_\beta=\frac{1}{N+\sigma_{0,\mathrm{old}}^{-2}},
\]
then
\[
\boxed{
\sigma_{0,\mathrm{new}}^2
=
\sum_{j=1}^J \pi_j\left(m_{\beta j}^2+v_\beta\right).
}
\]

For $\nu_0$, the same EM view reduces the update to a one-dimensional
optimization, but not to a closed-form expression. Let $p=J-1$ and write the
current posterior of $S_j$ as
\[
S_j\mid x,\gamma=j
\sim
\mathcal{IW}_p(\Psi_j,m_j),
\qquad
m_j=\nu_0^{\mathrm{old}}+J+2,
\]
where
\[
\Psi_j
=
\nu_0^{\mathrm{old}}\bar S_j+
\frac{N\nu_0^{\mathrm{old}}}{\nu_0^{\mathrm{old}}+Nx_j^2}
(x_{-j}-x_j\mu_j)(x_{-j}-x_j\mu_j)^\top.
\]
Also let
\[
\bar z_j
=
\frac{\nu_0^{\mathrm{old}}\mu_j+Nx_jx_{-j}}
{\nu_0^{\mathrm{old}}+Nx_j^2},
\qquad
\kappa_j^{\mathrm{old}}=\nu_0^{\mathrm{old}}+Nx_j^2.
\]
The needed posterior expectations are
\[
A_j=\E_q\{\tr(\bar S_jS_j^{-1})\}
=
m_j\tr(\bar S_j\Psi_j^{-1}),
\]
\[
B_j=\E_q\{(z_j-\mu_j)^\top S_j^{-1}(z_j-\mu_j)\}
=
m_j(\bar z_j-\mu_j)^\top\Psi_j^{-1}(\bar z_j-\mu_j)
+\frac{p}{\kappa_j^{\mathrm{old}}},
\]
and
\[
C_j=\E_q\{\log|S_j|\}
=
\log|\Psi_j|
-\sum_{r=1}^p \psi\left(\frac{m_j+1-r}{2}\right)
-p\log 2.
\]
Up to constants that do not depend on $\nu_0$, the EM objective for $\nu_0$ is
\begin{align*}
Q_\nu(\nu_0)
=
\sum_{j=1}^J \pi_j
\Bigg[
&\frac{p}{2}\log\nu_0
+\frac{\nu_0+J+1}{2}\log|\nu_0\bar S_j|
-\frac{(\nu_0+J+1)p}{2}\log 2 \\
&-\log\Gamma_p\left(\frac{\nu_0+J+1}{2}\right)
-\frac{\nu_0+2J+2}{2}C_j
-\frac{\nu_0}{2}(A_j+B_j)
\Bigg].
\end{align*}
Thus
\[
\boxed{
\nu_{0,\mathrm{new}}=\arg\max_{\nu_0>0} Q_\nu(\nu_0).
}
\]
So the exact-posterior ELBO gives a particularly nice update for
$\sigma_0^2$, and a stable scalar optimization for $\nu_0$. For $\nu_0$, the
closed-form marginal likelihood $\mathcal L(\nu_0,\sigma_0^2)$ above is often
even simpler to optimize directly, but the EM form is useful for deriving a
monotone coordinate-ascent algorithm.

\section{Extension to susie}
Now we try to extend model B into a Susie-like model, where we allow more than one non-zero effect.

\end{document}
